% Pursuing Stacks (cotes 134-2 to 134-8): archive cote 134-7, site page 62 (PDF page 63), typescript page 504.
% typescript. Draft compared: ps6.tex:3039-3098.
% AI-transcribed by Claude Opus 5.5 (claude-opus-5-5, adaptive thinking, effort high), Message Batches API, batch ps4 (msgbatch_012Aoc1NuskATPsyBeKSYK7J), 2026-11-12; not yet checked by eye.
\scanpage{134-7/62}
\mspage{504}

pursuing even for the \inserted{mere} sake of studying ordinary homotopy types.  The main
reason for feeling this way is the amazingly simple description of the
homotopy modules $\hand{\pi}_{i}$ of a homotopy type defined in terms of a semisimplicial
(or hemispherical) unipotent bundle,  \xout{which} as derived functors \xout{of the}
(so to say) of the Lie functor (cf. previous section 118).  The main test
for deciding whether there is indeed a rewarding new tool to be dug out,
is to see whether or not in the model categories \xout{\illegible{1}} $\hand{\mathcal{M}}_{n}(k)$ we have been
working with so far,  \xout{\illegible{1}} the standart homotopy constructions (around loop-%
spaces and suspensions) make a sense,  and in such a way of course that
the canonical functor from schematic to discrete homotopy types should com-%
mute to these operations.  It may well turn out that to get a handy forma-%
lism,  one will have still to modify more or less the conceptual set-up
of unipotent bundles I've been tentatively working with so far.  I already
\inserted{lately} hit upon suggestions of such modifications \deleted{lately}, and \unsure{presumably|presumahly}  I'm
going to discuss this still,  before leaving the topic of schematizations.

Another reason which makes me feel that there should exist a notion of
homotopy types over more general groundrings \inserted{$k$} than $\underline{Z}$ ,  is that for a number
of rings,  such a notion has been known for quite a while.  If I got it \xout{rig}
right,  already in the late sixties (even before I withdrew from the mathe-%
matical milieu) I heard about such things as\xout{s} homotopy types over residue\xout{s}
\inserted{class} rings $\underline{Z}/n\underline{Z}$ ,  or over rings(such as $\underline{Q}$) which are localizations of $\underline{Z}$ ,  or
over rings such as $\underline{Z}^{\wedge}$ \xout{\illegible{1}} or $\underline{Z}_{p}$ (p-adic integers)  which are completions
of $\underline{Z}$ with respect to a suitable  \xout{\unsure{padic topolgy}} linear topology.  Last
week,  which was the first time I was at the university \xout{since} after the
summer vacations,  I took from the library the Bousfield-Kan \inserted[typed]{Lecture Notes} book on
homotopy limits (which had been pointed out to me by Tim Porter in June,
when he had taken the trouble to tell me about``shape theory'' and it's rela-%
tions to \inserted[typed]{(filtering)} \deleted{homotopy} limits), and while glancing through it,  I noticed there
\xout{was} is a systematic treatment of such homotopy \uline{types.}\hand{$\rfloor\!\ulcorner$}

At this very minute
I had a closer look upon the introduction of part I,  it turns out that
\xout{they} Bousfield and Kan are working with an arbitrary (commutative) ground-%
ring k,  and they are defining corresponding k-completion  $k_{\infty}X$  of a homo-%
topy type k,  rather than a notion of ``homotopy type over k''.  But the two
kind of notions are surely closely related,  the k-completion of BK presu-%
mably should have more or less the meaning of ground \uline{ring extension} $\underline{Z}\hand{\to}$ k. \margin{definitely not, in general!}
At any rate,  for 1-connected spaces and k=$\underline{Z}$  the completion operation
seems to be no more no less than just the identity,  thus it would seem
that the implicit notion of ``homotopy type over $\underline{Z}$ ''  should be just the
ordinary ``discrete'' notion of homotopy types  -  unlike the notion of

\margin{62}
